\chapter{Benchmarking ST against GD32}
\label{chap:vs-gd32}

\section*{Introduction}

GigaDevice Semiconductor is a Chinese fabless semiconductor company that produces the GD32
series of 32-bit microcontrollers \cite{gigadevice-gd32}. The GD32 family is often
positioned as a pin-compatible, drop-in alternative to STM32 parts, sharing the same Arm
Cortex-M cores, similar peripheral sets, and comparable pricing. This close correspondence
makes GD32 a particularly relevant competitor: engineers evaluating an MCU platform
frequently shortlist both vendors, and the choice often comes down to measurable differences
in software-stack quality rather than silicon specifications.

In this chapter we compare the STM32Cube ecosystem against GigaDevice's GD32 firmware
libraries on two middleware stacks: the \textbf{USB device stack} (exercised through HID and
CDC scenarios) and the \textbf{FreeRTOS} real-time operating system integration (exercised
through basic, cooperative, and preemptive scheduling scenarios). For each stack we present
the scenario design, per-scenario footprint and execution-time results, a per-layer ROM
breakdown, and an interpretation of the trade-offs the data reveals.

All measurements follow the methodology described in
Chapter~\ref{chap:general-description}: identical scenarios on both platforms, same compiler
optimisation level, same clock frequency, and results extracted from IAR EWARM linker maps
using the footprint analysis tool presented in Chapter~\ref{chap:footprint-tool}.

% ====================================================================
\section{Hardware platforms}
\label{sec:gd32-boards}

Every scenario in this chapter runs on two boards hosting comparable Cortex-M33 parts: the
STMicroelectronics \textbf{NUCLEO-U575ZI-Q} (STM32U575ZIT6Q) and GigaDevice's
\textbf{GD32E507} evaluation board (GD32E507ZET6). Both devices integrate an Arm Cortex-M33
core clocked at 160\,MHz for these tests, ensuring that performance differences reflect
software-stack efficiency rather than raw silicon speed.
Table~\ref{tab:gd32-boards} contrasts their key characteristics.

\begin{table}[H]
  \centering
  \caption{Board characteristics: ST reference vs GigaDevice GD32.}
  \label{tab:gd32-boards}
  \begin{tabularx}{\textwidth}{l X X}
    \toprule
    \textbf{Characteristic} & \textbf{ST (STM32U575ZIT6Q)} & \textbf{GD32 (GD32E507ZET6)} \\
    \midrule
    Core              & Arm Cortex-M33              & Arm Cortex-M33 \\
    Clock (test)      & 160\,MHz                    & 160\,MHz \\
    Flash             & 2\,MB                       & 512\,KB \\
    SRAM              & 786\,KB                     & 128\,KB \\
    USB               & Full-Speed device            & Full-Speed device \\
    Security          & TrustZone, AES, PKA, HASH    & --- \\
    Toolchain / IDE   & STM32CubeIDE / IAR EWARM     & IAR EWARM \\
    % TODO(source): verify GD32E507 reference price
    \bottomrule
  \end{tabularx}
\end{table}

The STM32U575 belongs to the ultra-low-power STM32U5 series and offers substantially more
flash and RAM, hardware security accelerators, and a richer peripheral set
\cite{stm32u575-ds}. The GD32E507 targets a similar performance point at a competitive
price but with fewer on-chip resources \cite{gigadevice-gd32}. For our benchmarks, the
relevant constant is that \emph{both devices run the same core at the same frequency},
so footprint and execution-time differences are attributable to the vendor's software stack
and HAL implementation.

% ====================================================================
% STACK 1: USB DEVICE STACK
% ====================================================================
\section{USB device stack}
\label{sec:gd32-usb}

The Universal Serial Bus (USB) device stack is a critical middleware component for any MCU
that needs to communicate with a host PC. We benchmark ST's USB Device Library (part of the
STM32Cube middleware) against GigaDevice's USB library, both operating in
\textbf{Full-Speed} mode and running bare-metal (no RTOS).

% --- Scenario design ------------------------------------------------
\subsection{Scenario design}
\label{sec:gd32-usb-scenarios}

We defined two scenarios that cover the most commonly used USB device classes:

\begin{enumerate}
  \item \textbf{HID (Human Interface Device)} --- a low-rate interrupt-transfer scenario in
    which the MCU acts as a USB mouse, sending periodic HID reports to the host. This
    exercises the enumeration path, descriptor handling, and interrupt-endpoint transfer
    logic.
  \item \textbf{CDC (Communication Device Class)} --- a higher-throughput bulk-transfer
    scenario in which the MCU exposes a virtual COM port. The host sends data to the
    device, which echoes it back. This exercises bulk IN/OUT transfers, line-coding
    negotiation, and buffering.
\end{enumerate}

Each scenario was implemented identically on both platforms: same USB descriptors, same
endpoint configuration (Full-Speed), same clock source (internal oscillator ---
HSI48 on STM32, IRC48M on GD32), and same core frequency of 160\,MHz.
Table~\ref{tab:gd32-usb-scenarios} summarises the scenarios and the metrics collected.

\begin{table}[H]
  \centering
  \caption{USB device stack benchmark scenarios.}
  \label{tab:gd32-usb-scenarios}
  \begin{tabularx}{\textwidth}{l X l}
    \toprule
    \textbf{Scenario} & \textbf{What it exercises} & \textbf{Metrics} \\
    \midrule
    HID  & Interrupt transfers, enumeration, descriptor handling
         & ROM/RAM (bytes), init/enum/TX time (ms) \\
    CDC  & Bulk transfers, virtual COM port, echo loopback
         & ROM/RAM (bytes), init/enum/TX/RX time (ms) \\
    \bottomrule
  \end{tabularx}
\end{table}

% --- Results --------------------------------------------------------
\subsection{Results}
\label{sec:gd32-usb-results}

% ---- HID ----------------------------------------------------------
\subsubsection{HID scenario}
\label{sec:gd32-usb-hid}

Table~\ref{tab:gd32-usb-hid-footprint} presents the total ROM and RAM footprint for the
HID scenario.

\begin{table}[H]
  \centering
  \caption{USB HID device footprint (bytes).}
  \label{tab:gd32-usb-hid-footprint}
  \begin{tabular}{l r r r}
    \toprule
    \textbf{Metric} & \textbf{STM32U575} & \textbf{GD32E507} & \textbf{Difference} \\
    \midrule
    ROM (total) & 16\,269 & 9\,175  & $+77\,\%$ \\
    RAM (total) & 3\,446  & 9\,608  & $-65\,\%$ \\
    \bottomrule
  \end{tabular}
\end{table}

The STM32 firmware consumes 77\,\% more ROM than GD32, but uses 65\,\% \emph{less} RAM.
The RAM advantage is substantial: the GD32 implementation requires 0x2000 bytes (8\,KB) of
stack and 0x2000 bytes of heap, whereas the STM32 configuration needs only 0x400 bytes
(1\,KB) of stack and 0x200 bytes (512~bytes) of heap.

Table~\ref{tab:gd32-usb-hid-rom} breaks down the ROM consumption by architectural layer.

\begin{table}[H]
  \centering
  \caption{USB HID ROM breakdown by layer (bytes).}
  \label{tab:gd32-usb-hid-rom}
  \begin{tabular}{l r r r}
    \toprule
    \textbf{Layer} & \textbf{STM32U575} & \textbf{GD32E507} & \textbf{Difference} \\
    \midrule
    Application    & 2\,580  & 1\,942 & $+33\,\%$ \\
    HAL drivers    & 10\,426 & 820    & $+1\,171\,\%$ \\
    MW stack       & 2\,859  & 5\,743 & $-50\,\%$ \\
    \bottomrule
  \end{tabular}
\end{table}

The HAL layer dominates the difference. On STM32, the HAL RCC module alone accounts for
4\,560 bytes (44\,\% of total HAL ROM) and the combined HAL/LL USB driver adds another
4\,988 bytes (48\,\%). By contrast, GD32's HAL layer is minimal at 820 bytes because
GigaDevice folds much of the peripheral-initialisation logic into the middleware itself.
This architectural choice gives GD32 a smaller HAL at the cost of a larger middleware
layer: the GD32 USB library ROM (5\,743 bytes) is twice the size of ST's (2\,859 bytes).

The STM32 application layer is 33\,\% larger because the STM32Cube USB library externalises
device configuration into user files (\texttt{usbd\_desc.c}, \texttt{usbd\_conf.c}),
accounting for 985 bytes (38\,\% of application ROM). On GD32, these configuration concerns
are handled inside the middleware. This design gives STM32 users more configuration
flexibility at the cost of a larger application footprint.

Table~\ref{tab:gd32-usb-hid-exec} shows the execution-time measurements.

\begin{table}[H]
  \centering
  \caption{USB HID execution time (ms).}
  \label{tab:gd32-usb-hid-exec}
  \begin{tabular}{l r r r}
    \toprule
    \textbf{Operation} & \textbf{STM32U575} & \textbf{GD32E507} & \textbf{Difference} \\
    \midrule
    Initialisation      & 9.990  & 26.084 & $-61.7\,\%$ \\
    Enumeration         & 213    & 216    & $-1.4\,\%$ \\
    Data transmission   & 0.0093 & 0.0094 & $-0.7\,\%$ \\
    \bottomrule
  \end{tabular}
\end{table}

The STM32 USB initialisation is 61.7\,\% faster than GD32's --- a significant difference
that reflects more efficient clock-tree configuration and peripheral bring-up in the
STM32Cube HAL. Enumeration and data-transmission times are nearly identical, as expected:
these phases are dominated by USB protocol timing rather than software overhead.

% ---- CDC ----------------------------------------------------------
\subsubsection{CDC scenario}
\label{sec:gd32-usb-cdc}

Table~\ref{tab:gd32-usb-cdc-footprint} presents the total footprint for the CDC scenario.

\begin{table}[H]
  \centering
  \caption{USB CDC device footprint (bytes).}
  \label{tab:gd32-usb-cdc-footprint}
  \begin{tabular}{l r r r}
    \toprule
    \textbf{Metric} & \textbf{STM32U575} & \textbf{GD32E507} & \textbf{Difference} \\
    \midrule
    ROM (total) & 17\,049 & 9\,085  & $+88\,\%$ \\
    RAM (total) & 4\,089  & 9\,683  & $-57.7\,\%$ \\
    \bottomrule
  \end{tabular}
\end{table}

The pattern mirrors HID: STM32 uses substantially more ROM ($+88\,\%$) but far less RAM
($-57.7\,\%$). The CDC scenario adds transmit and receive buffer management, which
increases both vendors' footprints relative to HID, but the ratio remains consistent.

Table~\ref{tab:gd32-usb-cdc-rom} provides the per-layer ROM breakdown.

\begin{table}[H]
  \centering
  \caption{USB CDC ROM breakdown by layer (bytes).}
  \label{tab:gd32-usb-cdc-rom}
  \begin{tabular}{l r r r}
    \toprule
    \textbf{Layer} & \textbf{STM32U575} & \textbf{GD32E507} & \textbf{Difference} \\
    \midrule
    Application    & 2\,683  & 1\,882 & $+42.5\,\%$ \\
    HAL drivers    & 10\,446 & 820    & $+1\,174\,\%$ \\
    MW stack       & 2\,423  & 5\,709 & $-57.5\,\%$ \\
    \bottomrule
  \end{tabular}
\end{table}

The root cause is the same as in HID. The STM32 HAL RCC module (4\,560 bytes) and the
HAL/LL USB driver (5\,006 bytes) together account for 92\,\% of the HAL ROM.
The STM32 application layer is 42.5\,\% larger due to the externalised USB device
configuration files (\texttt{usb\_device.c}, \texttt{usbd\_cdc\_if.c},
\texttt{usbd\_desc.c}, \texttt{usbd\_conf.c}), contributing 1\,164 bytes (43\,\% of
application ROM). In return, the STM32 USB middleware library itself is 57.5\,\% smaller
than GD32's, because GigaDevice bundles configuration and driver logic into its middleware
layer.

Table~\ref{tab:gd32-usb-cdc-exec} shows the execution-time measurements.

\begin{table}[H]
  \centering
  \caption{USB CDC execution time (ms).}
  \label{tab:gd32-usb-cdc-exec}
  \begin{tabular}{l r r r}
    \toprule
    \textbf{Operation} & \textbf{STM32U575} & \textbf{GD32E507} & \textbf{Difference} \\
    \midrule
    Initialisation      & 9.994   & 26.076 & $-61.7\,\%$ \\
    Enumeration         & 218.2   & 218.5  & $-0.1\,\%$ \\
    Data transmission   & 0.00988 & 0.01003 & $-1.5\,\%$ \\
    Data reception      & 0.00986 & 0.00993 & $-0.7\,\%$ \\
    \bottomrule
  \end{tabular}
\end{table}

Once again, initialisation is where STM32 dominates: 61.7\,\% faster than GD32. The
remaining operations (enumeration, transmission, reception) are within 1.5\,\% of each
other, confirming that steady-state USB throughput is protocol-bound rather than
software-bound.

% --- Interpretation -------------------------------------------------
\subsection{Interpretation}
\label{sec:gd32-usb-interp}

The USB benchmark reveals a clear architectural trade-off between the two vendors.

\paragraph{ROM footprint.}
STM32 consumes 77--88\,\% more total ROM than GD32 across both scenarios. The dominant
contributor is the HAL layer, which is roughly twelve times larger on STM32. Two modules
drive this cost: the \textbf{HAL RCC} clock-configuration driver ($\approx$4\,560 bytes)
and the \textbf{HAL/LL USB} peripheral driver ($\approx$5\,000 bytes). GigaDevice achieves
a much smaller HAL by embedding peripheral-initialisation logic inside the middleware
library rather than exposing it as a separate, reusable driver layer.

\paragraph{RAM footprint.}
The situation is reversed: STM32 uses 57--65\,\% less RAM. The GD32 implementation
defaults to much larger stack and heap allocations (8\,KB each vs $\leq$1.5\,KB for STM32).
This suggests that the GD32 USB library relies more heavily on dynamic allocation or deep
call chains, whereas the STM32Cube library uses statically allocated buffers and a
shallower call stack.

\paragraph{Execution time.}
STM32's USB initialisation is consistently 61.7\,\% faster across both HID and CDC. After
initialisation, the two platforms perform within 1--2\,\% of each other, as the USB
protocol itself (enumeration handshake, bulk/interrupt transfers) dictates the timing.

\paragraph{Design philosophy.}
STM32Cube separates concerns: HAL handles clocks and peripheral access, the middleware
handles protocol logic, and the application owns configuration through user-editable files
(\texttt{usbd\_conf.c}, \texttt{usbd\_desc.c}). This modular design gives developers
explicit control and reusability at the cost of a larger total ROM. GigaDevice takes a more
monolithic approach --- smaller HAL, but a thicker middleware that bundles configuration
internally, reducing ROM but also reducing the user's configuration flexibility.

\paragraph{Enhancement proposals.}
The data suggests two concrete areas where STM32Cube could reduce ROM without sacrificing
functionality:

\begin{enumerate}
  \item \textbf{HAL RCC optimisation.} The HAL RCC module accounts for $\approx$4.5\,KB in
    every USB project regardless of the clocks actually configured. A conditionally compiled
    or link-time-optimised RCC driver that includes only the clock paths used would
    significantly narrow the HAL gap.
  \item \textbf{HAL/LL USB driver consolidation.} The combined HAL and LL USB driver
    ($\approx$5\,KB) provides a comprehensive abstraction but may include code paths
    (e.g.\ High-Speed, OTG, DMA) that Full-Speed-only projects never execute. Finer
    compile-time selection of USB features could reduce this cost.
\end{enumerate}

% ====================================================================
% STACK 2: FreeRTOS
% ====================================================================
\section{FreeRTOS integration}
\label{sec:gd32-freertos}

FreeRTOS is the most widely adopted real-time operating system for Cortex-M
microcontrollers \cite{freertos}. Both STMicroelectronics and GigaDevice ship FreeRTOS
integration packages as part of their firmware offerings. We benchmark the two
implementations using FreeRTOS v11.2, running on the same two boards at 160\,MHz with a
1\,ms tick (time slice) per task.

% --- Scenario design ------------------------------------------------
\subsection{Scenario design}
\label{sec:gd32-freertos-scenarios}

We defined three scheduling scenarios that stress progressively more complex RTOS
mechanisms:

\begin{enumerate}
  \item \textbf{Basic processing} --- a single counting thread runs at low priority
    alongside a high-priority reporting thread. The counting thread increments a global
    counter in a tight loop; the reporting thread wakes every 30\,s and prints the counter
    delta. This isolates single-task throughput with minimal scheduler overhead.
  \item \textbf{Cooperative processing} --- five threads at the same low priority, each
    incrementing its own counter and then calling \texttt{taskYIELD()} to voluntarily yield
    the processor. A high-priority reporter prints all five counters every 30\,s. This
    exercises voluntary context switching under equal-priority round-robin scheduling.
  \item \textbf{Preemptive processing} --- five threads at ascending priorities. Each
    thread resumes the next-higher-priority thread (which preempts it), increments its
    counter, then suspends itself. A high-priority reporter prints all five counters every
    30\,s. This exercises the full preemptive scheduler: suspend, resume, and
    priority-based preemption.
\end{enumerate}

Table~\ref{tab:gd32-freertos-scenarios} summarises the three scenarios.

\begin{table}[H]
  \centering
  \caption{FreeRTOS benchmark scenarios.}
  \label{tab:gd32-freertos-scenarios}
  \begin{tabularx}{\textwidth}{l X l}
    \toprule
    \textbf{Scenario} & \textbf{What it exercises} & \textbf{Metric} \\
    \midrule
    Basic        & Single-thread throughput, minimal scheduler overhead
                 & Thread iterations/s \\
    Cooperative  & Voluntary context switching (\texttt{taskYIELD})
                 & Total thread iterations/30\,s \\
    Preemptive   & Priority-based preemption, suspend/resume
                 & Total thread iterations/30\,s \\
    \bottomrule
  \end{tabularx}
\end{table}

% --- Results --------------------------------------------------------
\subsection{Results}
\label{sec:gd32-freertos-results}

% ---- Basic processing ---------------------------------------------
\subsubsection{Basic processing}
\label{sec:gd32-freertos-basic}

Table~\ref{tab:gd32-freertos-basic-perf} shows the performance result and
Table~\ref{tab:gd32-freertos-basic-fp} the footprint.

\begin{table}[H]
  \centering
  \caption{FreeRTOS basic processing --- performance.}
  \label{tab:gd32-freertos-basic-perf}
  \begin{tabular}{l r r r}
    \toprule
    \textbf{Metric} & \textbf{STM32U575} & \textbf{GD32E507} & \textbf{Difference} \\
    \midrule
    Thread iterations/s & 47\,943 & 47\,955 & $-0.025\,\%$ \\
    \bottomrule
  \end{tabular}
\end{table}

\begin{table}[H]
  \centering
  \caption{FreeRTOS basic processing --- footprint (bytes).}
  \label{tab:gd32-freertos-basic-fp}
  \begin{tabular}{l r r r}
    \toprule
    \textbf{Metric} & \textbf{STM32U575} & \textbf{GD32E507} & \textbf{Difference} \\
    \midrule
    ROM & 14\,266 & 10\,128 & $+40.9\,\%$ \\
    RAM & 11\,960 & 11\,716 & $+2.1\,\%$ \\
    \bottomrule
  \end{tabular}
\end{table}

Performance is virtually identical ($-0.025\,\%$), confirming that the counting loop is
CPU-bound and both FreeRTOS ports schedule it with the same efficiency. The ROM difference
($+40.9\,\%$) is significant; the RAM difference is negligible.

% ---- Cooperative processing ----------------------------------------
\subsubsection{Cooperative processing}
\label{sec:gd32-freertos-coop}

\begin{table}[H]
  \centering
  \caption{FreeRTOS cooperative processing --- performance.}
  \label{tab:gd32-freertos-coop-perf}
  \begin{tabular}{l r r r}
    \toprule
    \textbf{Metric} & \textbf{STM32U575} & \textbf{GD32E507} & \textbf{Difference} \\
    \midrule
    Thread iterations/30\,s & 1\,220\,066 & 1\,200\,666 & $+1.6\,\%$ \\
    \bottomrule
  \end{tabular}
\end{table}

\begin{table}[H]
  \centering
  \caption{FreeRTOS cooperative processing --- footprint (bytes).}
  \label{tab:gd32-freertos-coop-fp}
  \begin{tabular}{l r r r}
    \toprule
    \textbf{Metric} & \textbf{STM32U575} & \textbf{GD32E507} & \textbf{Difference} \\
    \midrule
    ROM & 15\,422 & 11\,224 & $+37\,\%$ \\
    RAM & 10\,972 & 10\,728 & $+2.3\,\%$ \\
    \bottomrule
  \end{tabular}
\end{table}

STM32 completes 1.6\,\% more thread iterations over 30 seconds, indicating a marginally
faster \texttt{taskYIELD()} context switch. The footprint gap narrows slightly to
$+37\,\%$ ROM.

% ---- Preemptive processing -----------------------------------------
\subsubsection{Preemptive processing}
\label{sec:gd32-freertos-preemptive}

\begin{table}[H]
  \centering
  \caption{FreeRTOS preemptive processing --- performance.}
  \label{tab:gd32-freertos-preemptive-perf}
  \begin{tabular}{l r r r}
    \toprule
    \textbf{Metric} & \textbf{STM32U575} & \textbf{GD32E507} & \textbf{Difference} \\
    \midrule
    Thread iterations/30\,s & 294\,615 & 285\,144 & $+3.3\,\%$ \\
    \bottomrule
  \end{tabular}
\end{table}

\begin{table}[H]
  \centering
  \caption{FreeRTOS preemptive processing --- footprint (bytes).}
  \label{tab:gd32-freertos-preemptive-fp}
  \begin{tabular}{l r r r}
    \toprule
    \textbf{Metric} & \textbf{STM32U575} & \textbf{GD32E507} & \textbf{Difference} \\
    \midrule
    ROM & 16\,226 & 12\,028 & $+35\,\%$ \\
    RAM & 10\,996 & 10\,748 & $+2.3\,\%$ \\
    \bottomrule
  \end{tabular}
\end{table}

Preemptive scheduling is the most demanding scenario, involving suspend, resume, and
priority-based task switching. STM32 achieves 3.3\,\% more iterations, the largest
performance advantage across the three FreeRTOS scenarios. The ROM gap ($+35\,\%$) is
consistent with the other scenarios.

% ---- Per-layer ROM breakdown (all three scenarios) -----------------
\subsubsection{Per-layer ROM breakdown}
\label{sec:gd32-freertos-rom-breakdown}

To understand \emph{where} the ROM overhead originates, Table~\ref{tab:gd32-freertos-rom}
breaks down the ROM footprint into three architectural layers --- application, HAL drivers,
and FreeRTOS middleware stack --- for all three scenarios.

\begin{table}[H]
  \centering
  \caption{FreeRTOS ROM breakdown by layer (bytes).}
  \label{tab:gd32-freertos-rom}
  \begin{tabular}{l l r r r}
    \toprule
    \textbf{Scenario} & \textbf{Layer} & \textbf{STM32U575} & \textbf{GD32E507}
        & \textbf{Diff.} \\
    \midrule
    \multirow{3}{*}{Basic}
      & Application & 4\,531  & 4\,278 & $+6\,\%$ \\
      & HAL drivers & 4\,927  & 1\,078 & $+357\,\%$ \\
      & MW stack    & 4\,808  & 4\,772 & $+0.7\,\%$ \\
    \midrule
    \multirow{3}{*}{Cooperative}
      & Application & 5\,823  & 5\,484 & $+6\,\%$ \\
      & HAL drivers & 4\,911  & 1\,078 & $+355\,\%$ \\
      & MW stack    & 4\,688  & 4\,662 & $+0.6\,\%$ \\
    \midrule
    \multirow{3}{*}{Preemptive}
      & Application & 6\,191  & 5\,810 & $+6\,\%$ \\
      & HAL drivers & 4\,911  & 1\,078 & $+355\,\%$ \\
      & MW stack    & 5\,124  & 5\,140 & $-0.3\,\%$ \\
    \bottomrule
  \end{tabular}
\end{table}

The pattern is strikingly consistent across all three scenarios:

\begin{itemize}
  \item The \textbf{FreeRTOS middleware stack} itself is virtually identical in size
    ($\leq$1\,\% difference). Both vendors ship the same upstream FreeRTOS kernel, so
    this is expected.
  \item The \textbf{application layer} is $\approx$6\,\% larger on STM32, reflecting
    slightly more initialisation code in the STM32Cube project template.
  \item The \textbf{HAL layer} is 355--357\,\% larger on STM32, entirely accounting for
    the total ROM gap. The dominant contributor is the \textbf{HAL RCC} module, which
    configures the clock tree.
\end{itemize}

% --- Interpretation -------------------------------------------------
\subsection{Interpretation}
\label{sec:gd32-freertos-interp}

\paragraph{Performance.}
FreeRTOS scheduling performance is nearly identical between the two platforms in the basic
scenario ($-0.025\,\%$) and slightly favours STM32 in the cooperative ($+1.6\,\%$) and
preemptive ($+3.3\,\%$) scenarios. The advantage grows with scheduler complexity, suggesting
that the STM32Cube FreeRTOS port has a marginally more efficient context-switch
implementation. However, these differences are small enough that application-level
performance is unlikely to be affected.

\paragraph{ROM footprint.}
The ROM overhead follows the same pattern as the USB stack: a 35--41\,\% total increase
driven almost entirely by the HAL layer, specifically the HAL RCC clock-configuration
module. Crucially, the FreeRTOS middleware itself contributes no measurable overhead ---
the gap is in the vendor's hardware-abstraction code, not in the RTOS kernel.

\paragraph{RAM footprint.}
RAM differences are negligible (2.1--2.3\,\%), indicating that both FreeRTOS ports allocate
task stacks and kernel structures similarly.

\paragraph{Enhancement proposals.}
The FreeRTOS results reinforce the conclusion from the USB benchmark: the HAL RCC module
is the single largest contributor to STM32's ROM overhead across \emph{all} middleware
stacks. A targeted effort to slim the RCC driver --- through dead-code elimination,
conditional compilation based on the clocks actually used, or link-time optimisation ---
would benefit every STM32 project, not just FreeRTOS.

% ====================================================================
\section*{Conclusion}

The ST-versus-GD32 comparison across two middleware stacks (USB device and FreeRTOS) and
five scenarios reveals a consistent pattern:

\begin{itemize}
  \item \textbf{Performance:} STM32 matches or slightly exceeds GD32 in every scenario.
    USB initialisation is 61.7\,\% faster; FreeRTOS preemptive scheduling is 3.3\,\% faster;
    steady-state throughput is equivalent.
  \item \textbf{ROM footprint:} STM32 consumes 35--88\,\% more flash, with the HAL layer
    (principally HAL RCC and, for USB, HAL/LL USB) responsible for the bulk of the
    difference. The middleware stacks themselves are comparable in size.
  \item \textbf{RAM footprint:} STM32 uses significantly less RAM than GD32 in USB
    scenarios (57--65\,\% less) thanks to smaller default stack/heap allocations and more
    static memory management. FreeRTOS RAM is nearly identical.
  \item \textbf{Architectural trade-off:} STM32Cube's modular, layered design (separate
    HAL, middleware, and user-configuration files) provides greater flexibility and
    reusability but at a measurable ROM cost. GigaDevice's more monolithic approach yields
    smaller ROM by bundling driver and configuration logic into the middleware, at the
    expense of less user-visible modularity.
\end{itemize}

The single most impactful enhancement for STM32 would be \textbf{optimising the HAL RCC
module}: it appears in every project and consistently accounts for the largest share of the
HAL ROM gap. A secondary improvement for USB projects would be \textbf{finer-grained
compile-time selection} of USB driver features (Full-Speed vs High-Speed, device vs OTG)
to exclude unused code paths.

\pagebreak